% !TEX root = report.tex
\section{Introduction and Literature Review}

\subsection{protic ionic liquids}
- {\pka}s are based on proton lability in aqueous environments and aren't representative of the dynamics in bulk IL solution.

- Hammett method for determining \pka
$$
H_0=pK(I)_{aq}+\log([I]_s/[IH]_s)
$$

- \pka can also be determined by titration

- ILs can have lots of reacting species and can form lots of weird byproducts, so a ``pure'' IL is typically considered to be one which has $<1\%$ neutral species. Higher than this can be considered mixtures of IL and neutral species

- aqueous \pka differences ($pK_A(base)-pK_a(acid)$) can be suggestive of how complete the proton transfer is in a PIL

- PILs with a $\Delta pK_a>8$ have nearly ideal Walden behaviour. \cite{Yoshizawa2003}?

- Walden plot of ionicity? one metric of how ionic a salt is - $\log(\text{equivalent conductivity})$vs $\log(\frac{1}{\eta})$

- The strength of the PIL acid has more of an impact on the ionicity than the size, with low \pka values leading to higher ionicity - a measure of how ionic a salt is.

- Liquid state has high viscosity, then as the temperature decreases you hit the melting point ($T_m$) where the liquid turns into a `plastic crystal' which is an amorphous non-crystalline solid. we keep decreasing the temperature, we hit the glass transition temperature ($T_g$) , corresponding to a maximum heat capacity, where the plastic crystal can fully crystallise.
In ILs, as the temperature decreases, the viscosity will gradually decrease, but won't have a definite melting point, it will just become more and more viscous.

- the glass transition temperature ($T_g$) is a ``qualitative signature'' of the ion mobility in the IL (10.1002/chem.200400817)

- Unsurprisingly, more hydrogen bonding means higher viscosity, as does more vdW interactions

- ion size doesn't affect imidazolium ILs that much (10.1021/ic951325x)

- The ability for PILs to catalyse a reaction through a Br{\o}nsted acid activated pathways was strongly correlated to the acid strength of the IL pair \cite{Greaves2008a}
	-- The acidity seems to be of the IL pair and not of the acidic component alone
	-- when using the same cation, the basicity of the anion (based on the \pka) dictated the relative acidity, and thus the catalytic ability of the ion pair (tested with [bbim] and [Hbim] and \cite{Palimkar2003}
	-- Furthers that the $\Delta pK_a$ can be used as a relative indicator of the catalytic potential of a PIL

\subsection{cyclisation of aminchalcone}

\cite{Du} showed \il{HBIm}{CF3SO3} to be the best for amine addition to chalcone, however this was in an aqueous environment. This used the IL as a Br{\o}nsted acid\\

\cite{Zheng2013} showed that piperidine as a Lewis base with \ce{KOH} was exceptionally efficient at converting the aminochalcone to the cyclised product (>99\% yield r.t.)

I'm guessing that this was adding in to the ketone, withdrawing electron density from the $\alpha,\beta$ double bond, making it more electrophilic, then the \ce{KOH} came through and mopped up the mopped up the proton form the resulting \ce{RNH3+}\\

\cite{Bunce2011} synthesised the product and made the flavone using just conc. \ce{HCl}

Would be activation by protonating the carbonyl, or the $\alpha$ site. according to the paper, these are both valid pathway